%% kapa: exam=1 round=32 year=2021 period=1 form=A subject=civil_law subject_ko=민법 seq=19 number=19
%% source: Q-Net 감정평가사 기출문제 (cst00309) artlSeq=5208397 file=extracted/1008/5208397/A형/2021년도 제32회 감정평가사 1차 1교시 A형.pdf
%% conversion: kordoc 4.24.51 -> tools/gen1.py (auto); review=proofread-sonnet; flags=none
\begin{question}{19}
소멸시효의 기산점에 관한 설명으로 옳지 않은 것은? (다툼이 있으면 판례에 따름)
\begin{choices}
  \item 정지조건부 권리는 조건이 성취되지 않은 동안에는 소멸시효가 진행되지 않는다.
  \item 이행기한을 정하지 않은 채권은 채권자의 이행최고가 있은 날로부터 소멸시효가 진행한다.
  \item 채무불이행으로 인한 손해배상청구권은 채무불이행시로부터 소멸시효가 진행한다.
  \item 동시이행의 항변권이 붙은 채권은 그 이행기로부터 소멸시효가 진행한다.
  \item 무권대리인에 대한 상대방의 계약이행청구권이나 손해배상청구권은 그 선택권을 행사할 수 있을 때부터 소멸시효가 진행한다.
\end{choices}
\end{question}
\answer{2}
